\section{Empatizar: el problema y quién lo tiene}
\label{sec:empatizar}

\subsection{Descripción del problema}

\begin{description}[style=nextline,leftmargin=0pt,font=\color{verde}]
\item[¿Qué problema existe?]
Mucha gente quiere usar menos el celular y no lo logra sostener. Las herramientas que existen muestran estadísticas o bloquean apps, pero no le dan nada a la persona cuando lo consigue.
\item[¿Quién lo tiene?]
Adultos de 18 a 35 años que usan el teléfono todo el día y que ya probaron alguna app de tiempo de pantalla.
\item[¿En qué contexto ocurre?]
En la vida diaria: en casa, en el trabajo, en el colectivo. El celular siempre está a mano y las notificaciones compiten por la atención.
\item[¿Cómo se resuelve hoy?]
Con Bienestar digital de Android o Tiempo en pantalla de iOS, y con apps como Forest o one sec, que obligan a iniciar una sesión o agregan una espera antes de abrir otras apps.
\item[¿Qué dificultades tiene esa solución?]
Piden un esfuerzo activo (acordarse de iniciar la sesión, configurar límites), funcionan como castigo y el logro no tiene ningún efecto fuera del teléfono. Al tiempo se ignoran o se desinstalan.
\item[¿Por qué una app móvil mejoraría la situación?]
El dato que importa, cuánto tiempo estuvo el teléfono sin usarse, lo registra el propio sistema operativo. Solo una app instalada en el teléfono puede leerlo sin pedirle nada al usuario.
\end{description}

\subsection{Usuario principal}

\begin{table}[H]
\centering\small
\begin{tabularx}{\textwidth}{@{}l L@{}}
\toprule
\textbf{Atributo} & \textbf{Descripción} \\
\midrule
Características & Entre 18 y 35 años, usa el celular muchas horas por día, quiere bajar ese uso pero no lo sostiene. \\
Necesidad principal & Un motivo concreto para dejar el teléfono, no más información sobre cuánto lo usa. \\
Contexto de uso & Casi todo el uso es pasivo: la app trabaja en segundo plano. Se abre un rato para ver puntos o canjear. \\
Frecuencia & Entre 2 y 4 aperturas breves por día. Un canje por semana, aproximadamente. \\
Restricciones & Poco tiempo, una sola mano, a veces caminando o parado frente al mostrador de un local. \\
Conectividad & Intermitente. Durante la pausa el teléfono puede estar en modo avión o sin señal. \\
Condiciones situacionales & Luz variable al mostrar el código en un local; apuro porque hay gente esperando. \\
Dificultades actuales & Se olvida de activar las herramientas y pierde la motivación porque no recibe nada a cambio. \\
\bottomrule
\end{tabularx}
\end{table}

\begin{nota}[title={Persona: Martina, 27 años},coltitle=white,colbacktitle=verde,fonttitle=\bfseries]
Trabaja como diseñadora en Palermo y vive a diez cuadras de la oficina. Revisa el celular apenas se despierta y otra vez antes de dormir. Probó Forest dos semanas y lo dejó: se olvidaba de plantar el árbol. Le gustaría leer más y salir a caminar, y le encantan los cafés de especialidad del barrio.

\textit{“Sé que lo uso demasiado. Lo que no tengo es una razón para dejarlo.”}
\end{nota}

La persona es un arquetipo armado con supuestos del equipo. La idea es validarla con las entrevistas del plan de la sección~\ref{sec:testear}.

El comercio adherido es un usuario secundario. Publica premios y los entrega en el local, pero en esta versión no tiene una app propia: sus datos se cargan como datos semilla.

\subsection{Contexto de uso}

Seguimos la idea de computación ubicua de Mark Weiser: la mejor interacción con esta app es no tener que abrirla. El usuario nunca declara que empieza una pausa. Deja el teléfono y la app se ocupa del resto. Solo vuelve a aparecer para avisar que sumó puntos.

Hay dos momentos en los que sí hay interacción, y los dos son cortos:
\begin{itemize}
\item \textbf{Mirar cómo va}: abre la app desde la notificación, ve el saldo y cuánto le falta para su meta. Menos de 10 segundos.
\item \textbf{Canjear}: está en el local, con una mano, quizás con gente atrás. Necesita encontrar el premio, confirmar y mostrar un código que se lea bien.
\end{itemize}

\subsection{Recorrido actual del usuario}

Así se ve hoy el intento de usar menos el celular con las herramientas que existen. La última columna es lo que tomamos para el diseño.

\begin{table}[H]
\centering\small
\begin{tabularx}{\textwidth}{@{}l L l L@{}}
\toprule
\textbf{Etapa} & \textbf{Qué hace} & \textbf{Cómo se siente} & \textbf{Oportunidad} \\
\midrule
Decide    & “Esta semana lo uso menos.” & Motivada & Aprovechar el impulso sin pedir configuración \\
Configura & Instala una app, arma límites y horarios & Cansada & Cero configuración \\
Usa       & Tiene que acordarse de activar el modo foco & Se olvida & Detectar las pausas sola \\
Ve el resultado & Un gráfico o un árbol virtual & Indiferente & Un premio real y cerca \\
Abandona  & Desinstala o ignora los avisos & Culpa & Reglas amables, racha y meta \\
\bottomrule
\end{tabularx}
\end{table}
